\documentclass[aspectratio=169,11pt]{beamer}
\usetheme{Madrid}
\usecolortheme{default}
\setbeamertemplate{navigation symbols}{}
\setbeamertemplate{footline}[frame number]
\usepackage{enumitem}
\usepackage{amsmath,amssymb}
\usepackage{booktabs}
\usepackage{array}

\title{Databricks ML Associate}
\subtitle{Questions 301--400 (Fourth 100)}
\author{Practice Exam Deck}
\date{}

\begin{document}
	
	\begin{frame}
		\titlepage
	\end{frame}
	
	\begin{frame}{How to Use This Deck}
		\begin{itemize}
			\item 100 new questions (Q301--Q400), none repeated from Q1--Q300.
			\item Each question slide has four options (A--D).
			\item The next slide gives the correct answer and explains why the right option is right and why each wrong option is wrong.
			\item This set emphasizes areas not yet covered: CLI/REST API, cluster config, registry specifics, model flavors, DLT, profiling, time series, and governance.
		\end{itemize}
	\end{frame}
	
	% ---------------------------------------------------------------
	\begin{frame}{Question 301}
		Which tool lets you automate Databricks operations from a terminal?
		\begin{enumerate}[label=\Alph*.]
			\item Databricks CLI
			\item Databricks SQL Editor
			\item Notebook widgets
			\item MLflow UI
		\end{enumerate}
	\end{frame}
	\begin{frame}{Answer 301}
		\textbf{Correct: A} --- the Databricks CLI automates operations.\\
		\vspace{0.5em}
		\textbf{Wrong:}
		\begin{itemize}
			\item B is SQL editing.
			\item C is notebook UI.
			\item D is experiment tracking.
		\end{itemize}
	\end{frame}
	
	% ---------------------------------------------------------------
	\begin{frame}{Question 302}
		Which REST API endpoint lists clusters?
		\begin{enumerate}[label=\Alph*.]
			\item \texttt{GET /api/2.0/clusters/list}
			\item \texttt{GET /api/2.0/clusters/all}
			\item \texttt{GET /api/2.0/clusters/describe}
			\item \texttt{GET /api/2.0/clusters}
		\end{enumerate}
	\end{frame}
	\begin{frame}{Answer 302}
		\textbf{Correct: A} --- \texttt{/clusters/list} returns clusters.\\
		\vspace{0.5em}
		\textbf{Wrong:}
		\begin{itemize}
			\item B, C, D are not the correct endpoint.
		\end{itemize}
	\end{frame}
	
	% ---------------------------------------------------------------
	\begin{frame}{Question 303}
		Which authentication method is recommended for the Databricks CLI?
		\begin{enumerate}[label=\Alph*.]
			\item Personal access token or OAuth
			\item Username and password only
			\item SSH key only
			\item Anonymous access
		\end{enumerate}
	\end{frame}
	\begin{frame}{Answer 303}
		\textbf{Correct: A} --- tokens or OAuth are recommended.\\
		\vspace{0.5em}
		\textbf{Wrong:}
		\begin{itemize}
			\item B, C, D are not recommended.
		\end{itemize}
	\end{frame}
	
	% ---------------------------------------------------------------
	\begin{frame}{Question 304}
		What is a personal access token used for?
		\begin{enumerate}[label=\Alph*.]
			\item Authenticating API and CLI requests
			\item Training models
			\item Storing features
			\item Registering models
		\end{enumerate}
	\end{frame}
	\begin{frame}{Answer 304}
		\textbf{Correct: A} --- PATs authenticate requests.\\
		\vspace{0.5em}
		\textbf{Wrong:}
		\begin{itemize}
			\item B, C, D are unrelated.
		\end{itemize}
	\end{frame}
	
	% ---------------------------------------------------------------
	\begin{frame}{Question 305}
		Which component hosts the MLflow tracking backend in Databricks?
		\begin{enumerate}[label=\Alph*.]
			\item The Databricks-managed tracking server
			\item A local SQLite database only
			\item A notebook
			\item A cluster driver only
		\end{enumerate}
	\end{frame}
	\begin{frame}{Answer 305}
		\textbf{Correct: A} --- Databricks hosts a managed tracking server.\\
		\vspace{0.5em}
		\textbf{Wrong:}
		\begin{itemize}
			\item B, C, D are not the managed backend.
		\end{itemize}
	\end{frame}
	
	% ---------------------------------------------------------------
	\begin{frame}{Question 306}
		What is a cluster access mode?
		\begin{enumerate}[label=\Alph*.]
			\item It controls user isolation and supported workloads
			\item It sets the runtime version
			\item It sets the number of workers
			\item It sets the storage location
		\end{enumerate}
	\end{frame}
	\begin{frame}{Answer 306}
		\textbf{Correct: A} --- access modes define isolation and features.\\
		\vspace{0.5em}
		\textbf{Wrong:}
		\begin{itemize}
			\item B, C, D are other settings.
		\end{itemize}
	\end{frame}
	
	% ---------------------------------------------------------------
	\begin{frame}{Question 307}
		Which access mode is required for Unity Catalog-enabled clusters?
		\begin{enumerate}[label=\Alph*.]
			\item Shared or single-user with UC support
			\item No isolation
			\item Legacy only
			\item Personal only
		\end{enumerate}
	\end{frame}
	\begin{frame}{Answer 307}
		\textbf{Correct: A} --- UC requires supported access modes.\\
		\vspace{0.5em}
		\textbf{Wrong:}
		\begin{itemize}
			\item B, C, D are not the UC-required modes.
		\end{itemize}
	\end{frame}
	
	% ---------------------------------------------------------------
	\begin{frame}{Question 308}
		What does autoscaling do?
		\begin{enumerate}[label=\Alph*.]
			\item Adjusts the number of workers based on load
			\item Changes the runtime version
			\item Changes the access mode
			\item Changes the storage
		\end{enumerate}
	\end{frame}
	\begin{frame}{Answer 308}
		\textbf{Correct: A} --- autoscaling adds/removes workers.\\
		\vspace{0.5em}
		\textbf{Wrong:}
		\begin{itemize}
			\item B, C, D are unrelated.
		\end{itemize}
	\end{frame}
	
	% ---------------------------------------------------------------
	\begin{frame}{Question 309}
		What is a cluster policy?
		\begin{enumerate}[label=\Alph*.]
			\item A rule set that constrains cluster configuration
			\item A model version
			\item A feature table
			\item A notebook
		\end{enumerate}
	\end{frame}
	\begin{frame}{Answer 309}
		\textbf{Correct: A} --- policies constrain cluster settings.\\
		\vspace{0.5em}
		\textbf{Wrong:}
		\begin{itemize}
			\item B, C, D are unrelated.
		\end{itemize}
	\end{frame}
	
	% ---------------------------------------------------------------
	\begin{frame}{Question 310}
		What is a single-node cluster best suited for?
		\begin{enumerate}[label=\Alph*.]
			\item Small-scale development and testing
			\item Large-scale distributed training
			\item SQL BI only
			\item Streaming only
		\end{enumerate}
	\end{frame}
	\begin{frame}{Answer 310}
		\textbf{Correct: A} --- single-node is for small workloads.\\
		\vspace{0.5em}
		\textbf{Wrong:}
		\begin{itemize}
			\item B needs workers.
			\item C is SQL warehouses.
			\item D is not its main use.
		\end{itemize}
	\end{frame}
	
	% ---------------------------------------------------------------
	\begin{frame}{Question 311}
		What is a driver node responsible for in a cluster?
		\begin{enumerate}[label=\Alph*.]
			\item Coordinating work and running the Spark driver
			\item Storing all data
			\item Serving REST requests only
			\item Running SQL only
		\end{enumerate}
	\end{frame}
	\begin{frame}{Answer 311}
		\textbf{Correct: A} --- the driver coordinates execution.\\
		\vspace{0.5em}
		\textbf{Wrong:}
		\begin{itemize}
			\item B, C, D are not the driver's role.
		\end{itemize}
	\end{frame}
	
	% ---------------------------------------------------------------
	\begin{frame}{Question 312}
		What does a worker node do?
		\begin{enumerate}[label=\Alph*.]
			\item Executes tasks assigned by the driver
			\item Hosts the Spark driver
			\item Serves REST endpoints
			\item Stores the model registry
		\end{enumerate}
	\end{frame}
	\begin{frame}{Answer 312}
		\textbf{Correct: A} --- workers execute tasks.\\
		\vspace{0.5em}
		\textbf{Wrong:}
		\begin{itemize}
			\item B is the driver.
			\item C, D are other services.
		\end{itemize}
	\end{frame}
	
	% ---------------------------------------------------------------
	\begin{frame}{Question 313}
		Which runtime version should you select to use Delta Live Tables?
		\begin{enumerate}[label=\Alph*.]
			\item A DLT-compatible runtime version
			\item Any standard runtime
			\item A single-node runtime only
			\item A SQL warehouse
		\end{enumerate}
	\end{frame}
	\begin{frame}{Answer 313}
		\textbf{Correct: A} --- DLT requires a compatible runtime.\\
		\vspace{0.5em}
		\textbf{Wrong:}
		\begin{itemize}
			\item B, C, D are not DLT runtimes.
		\end{itemize}
	\end{frame}
	
	% ---------------------------------------------------------------
	\begin{frame}{Question 314}
		What is a Delta Live Tables pipeline?
		\begin{enumerate}[label=\Alph*.]
			\item A declarative pipeline for batch and streaming data
			\item A model
			\item A notebook
			\item A cluster
		\end{enumerate}
	\end{frame}
	\begin{frame}{Answer 314}
		\textbf{Correct: A} --- DLT builds declarative pipelines.\\
		\vspace{0.5em}
		\textbf{Wrong:}
		\begin{itemize}
			\item B, C, D are not pipelines.
		\end{itemize}
	\end{frame}
	
	% ---------------------------------------------------------------
	\begin{frame}{Question 315}
		What does DLT's \texttt{@dlt.table} decorator do?
		\begin{enumerate}[label=\Alph*.]
			\item Defines a table in a DLT pipeline
			\item Trains a model
			\item Registers a model
			\item Deploys an endpoint
		\end{enumerate}
	\end{frame}
	\begin{frame}{Answer 315}
		\textbf{Correct: A} --- it defines a DLT table.\\
		\vspace{0.5em}
		\textbf{Wrong:}
		\begin{itemize}
			\item B, C, D are unrelated.
		\end{itemize}
	\end{frame}
	
	% ---------------------------------------------------------------
	\begin{frame}{Question 316}
		What does DLT's \texttt{@dlt.view} decorator do?
		\begin{enumerate}[label=\Alph*.]
			\item Defines a temporary view in a DLT pipeline
			\item Defines a table
			\item Defines a model
			\item Defines a cluster
		\end{enumerate}
	\end{frame}
	\begin{frame}{Answer 316}
		\textbf{Correct: A} --- it defines a DLT view.\\
		\vspace{0.5em}
		\textbf{Wrong:}
		\begin{itemize}
			\item B, C, D are false.
		\end{itemize}
	\end{frame}
	
	% ---------------------------------------------------------------
	\begin{frame}{Question 317}
		What is a data profiling notebook used for?
		\begin{enumerate}[label=\Alph*.]
			\item Summarizing data quality and distributions
			\item Training models
			\item Deploying models
			\item Registering models
		\end{enumerate}
	\end{frame}
	\begin{frame}{Answer 317}
		\textbf{Correct: A} --- profiling summarizes data.\\
		\vspace{0.5em}
		\textbf{Wrong:}
		\begin{itemize}
			\item B, C, D are unrelated.
		\end{itemize}
	\end{frame}
	
	% ---------------------------------------------------------------
	\begin{frame}{Question 318}
		Which chart best shows the relationship between two numeric variables?
		\begin{enumerate}[label=\Alph*.]
			\item Scatter plot
			\item Histogram
			\item Bar chart
			\item Pie chart
		\end{enumerate}
	\end{frame}
	\begin{frame}{Answer 318}
		\textbf{Correct: A} --- scatter plots show two-variable relationships.\\
		\vspace{0.5em}
		\textbf{Wrong:}
		\begin{itemize}
			\item B is one variable's distribution.
			\item C is categorical.
			\item D is part-to-whole.
		\end{itemize}
	\end{frame}
	
	% ---------------------------------------------------------------
	\begin{frame}{Question 319}
		Which chart shows the distribution of one numeric variable?
		\begin{enumerate}[label=\Alph*.]
			\item Histogram
			\item Scatter plot
			\item Line chart
			\item Heatmap
		\end{enumerate}
	\end{frame}
	\begin{frame}{Answer 319}
		\textbf{Correct: A} --- histograms show distributions.\\
		\vspace{0.5em}
		\textbf{Wrong:}
		\begin{itemize}
			\item B is relationships.
			\item C is trends.
			\item D is a matrix view.
		\end{itemize}
	\end{frame}
	
	% ---------------------------------------------------------------
	\begin{frame}{Question 320}
		Which chart is best for trends over time?
		\begin{enumerate}[label=\Alph*.]
			\item Line chart
			\item Scatter plot
			\item Pie chart
			\item Bar chart
		\end{enumerate}
	\end{frame}
	\begin{frame}{Answer 320}
		\textbf{Correct: A} --- line charts show trends.\\
		\vspace{0.5em}
		\textbf{Wrong:}
		\begin{itemize}
			\item B, C, D are less suited to trends.
		\end{itemize}
	\end{frame}
	
	% ---------------------------------------------------------------
	\begin{frame}{Question 321}
		What is a correlation matrix used for?
		\begin{enumerate}[label=\Alph*.]
			\item Showing pairwise correlations between numeric features
			\item Training models
			\item Serving models
			\item Registering models
		\end{enumerate}
	\end{frame}
	\begin{frame}{Answer 321}
		\textbf{Correct: A} --- correlation matrices show pairwise relationships.\\
		\vspace{0.5em}
		\textbf{Wrong:}
		\begin{itemize}
			\item B, C, D are unrelated.
		\end{itemize}
	\end{frame}
	
	% ---------------------------------------------------------------
	\begin{frame}{Question 322}
		What is a boxplot useful for?
		\begin{enumerate}[label=\Alph*.]
			\item Showing quartiles and outliers
			\item Showing correlations
			\item Showing trends
			\item Showing class balance
		\end{enumerate}
	\end{frame}
	\begin{frame}{Answer 322}
		\textbf{Correct: A} --- boxplots show quartiles and outliers.\\
		\vspace{0.5em}
		\textbf{Wrong:}
		\begin{itemize}
			\item B, C, D are other visualizations.
		\end{itemize}
	\end{frame}
	
	% ---------------------------------------------------------------
	\begin{frame}{Question 323}
		Why is a holdout test set important after tuning?
		\begin{enumerate}[label=\Alph*.]
			\item It gives an unbiased final estimate
			\item It tunes hyperparameters
			\item It stores models
			\item It registers models
		\end{enumerate}
	\end{frame}
	\begin{frame}{Answer 323}
		\textbf{Correct: A} --- the test set gives a final estimate.\\
		\vspace{0.5em}
		\textbf{Wrong:}
		\begin{itemize}
			\item B tuning uses validation data.
			\item C, D are unrelated.
		\end{itemize}
	\end{frame}
	
	% ---------------------------------------------------------------
	\begin{frame}{Question 324}
		What is the risk of tuning on the test set?
		\begin{enumerate}[label=\Alph*.]
			\item The test set is no longer an unbiased estimate
			\item Training becomes slower
			\item Accuracy always drops
			\item The model cannot be saved
		\end{enumerate}
	\end{frame}
	\begin{frame}{Answer 324}
		\textbf{Correct: A} --- tuning on test data leaks it.\\
		\vspace{0.5em}
		\textbf{Wrong:}
		\begin{itemize}
			\item B, C, D are not the main risk.
		\end{itemize}
	\end{frame}
	
	% ---------------------------------------------------------------
	\begin{frame}{Question 325}
		What is nested cross-validation used for?
		\begin{enumerate}[label=\Alph*.]
			\item Tuning and estimating performance without optimistic bias
			\item Training one model
			\item Deploying a model
			\item Registering a model
		\end{enumerate}
	\end{frame}
	\begin{frame}{Answer 325}
		\textbf{Correct: A} --- nested CV separates tuning from evaluation.\\
		\vspace{0.5em}
		\textbf{Wrong:}
		\begin{itemize}
			\item B, C, D are unrelated.
		\end{itemize}
	\end{frame}
	
	% ---------------------------------------------------------------
	\begin{frame}{Question 326}
		Which MLflow model flavor supports scikit-learn models?
		\begin{enumerate}[label=\Alph*.]
			\item \texttt{mlflow.sklearn}
			\item \texttt{mlflow.tensorflow}
			\item \texttt{mlflow.pytorch}
			\item \texttt{mlflow.onnx}
		\end{enumerate}
	\end{frame}
	\begin{frame}{Answer 326}
		\textbf{Correct: A} --- \texttt{sklearn} is the scikit-learn flavor.\\
		\vspace{0.5em}
		\textbf{Wrong:}
		\begin{itemize}
			\item B, C, D are other flavors.
		\end{itemize}
	\end{frame}
	
	% ---------------------------------------------------------------
	\begin{frame}{Question 327}
		Which MLflow model flavor supports Spark ML models?
		\begin{enumerate}[label=\Alph*.]
			\item \texttt{mlflow.spark}
			\item \texttt{mlflow.sklearn}
			\item \texttt{mlflow.pytorch}
			\item \texttt{mlflow.tensorflow}
		\end{enumerate}
	\end{frame}
	\begin{frame}{Answer 327}
		\textbf{Correct: A} --- \texttt{spark} is the Spark ML flavor.\\
		\vspace{0.5em}
		\textbf{Wrong:}
		\begin{itemize}
			\item B, C, D are other flavors.
		\end{itemize}
	\end{frame}
	
	% ---------------------------------------------------------------
	\begin{frame}{Question 328}
		Why is a custom \texttt{pyfunc} model useful?
		\begin{enumerate}[label=\Alph*.]
			\item It wraps arbitrary prediction logic behind a standard interface
			\item It trains faster
			\item It stores data
			\item It registers models
		\end{enumerate}
	\end{frame}
	\begin{frame}{Answer 328}
		\textbf{Correct: A} --- pyfunc standardizes custom logic.\\
		\vspace{0.5em}
		\textbf{Wrong:}
		\begin{itemize}
			\item B, C, D are false.
		\end{itemize}
	\end{frame}
	
	% ---------------------------------------------------------------
	\begin{frame}{Question 329}
		What method must a custom pyfunc model implement?
		\begin{enumerate}[label=\Alph*.]
			\item \texttt{predict()}
			\item \texttt{fit()}
			\item \texttt{transform()}
			\item \texttt{evaluate()}
		\end{enumerate}
	\end{frame}
	\begin{frame}{Answer 329}
		\textbf{Correct: A} --- pyfunc models implement \texttt{predict()}.\\
		\vspace{0.5em}
		\textbf{Wrong:}
		\begin{itemize}
			\item B, C, D are not required for prediction.
		\end{itemize}
	\end{frame}
	
	% ---------------------------------------------------------------
	\begin{frame}{Question 330}
		What does \texttt{mlflow.pyfunc.log\_model()} do?
		\begin{enumerate}[label=\Alph*.]
			\item Logs a custom pyfunc model as an MLflow artifact
			\item Registers a model
			\item Trains a model
			\item Deploys a model
		\end{enumerate}
	\end{frame}
	\begin{frame}{Answer 330}
		\textbf{Correct: A} --- it logs a pyfunc model.\\
		\vspace{0.5em}
		\textbf{Wrong:}
		\begin{itemize}
			\item B, C, D are other operations.
		\end{itemize}
	\end{frame}
	
	% ---------------------------------------------------------------
	\begin{frame}{Question 331}
		What does model serving require before an endpoint can be created?
		\begin{enumerate}[label=\Alph*.]
			\item A registered model version
			\item A feature table
			\item A notebook
			\item A cluster
		\end{enumerate}
	\end{frame}
	\begin{frame}{Answer 331}
		\textbf{Correct: A} --- serving needs a registered version.\\
		\vspace{0.5em}
		\textbf{Wrong:}
		\begin{itemize}
			\item B, C, D are not required.
		\end{itemize}
	\end{frame}
	
	% ---------------------------------------------------------------
	\begin{frame}{Question 332}
		What is a serving endpoint's scale-to-zero behavior?
		\begin{enumerate}[label=\Alph*.]
			\item It shuts down idle capacity to save cost
			\item It deletes the model
			\item It retrains the model
			\item It changes the schema
		\end{enumerate}
	\end{frame}
	\begin{frame}{Answer 332}
		\textbf{Correct: A} --- scale-to-zero saves cost when idle.\\
		\vspace{0.5em}
		\textbf{Wrong:}
		\begin{itemize}
			\item B, C, D are false.
		\end{itemize}
	\end{frame}
	
	% ---------------------------------------------------------------
	\begin{frame}{Question 333}
		What is a feature specification used for in online serving?
		\begin{enumerate}[label=\Alph*.]
			\item It tells the endpoint which online table to query for features
			\item It trains the model
			\item It registers the model
			\item It deploys the model
		\end{enumerate}
	\end{frame}
	\begin{frame}{Answer 333}
		\textbf{Correct: A} --- feature specs connect endpoints to online tables.\\
		\vspace{0.5em}
		\textbf{Wrong:}
		\begin{itemize}
			\item B, C, D are unrelated.
		\end{itemize}
	\end{frame}
	
	% ---------------------------------------------------------------
	\begin{frame}{Question 334}
		What is online feature serving latency primarily affected by?
		\begin{enumerate}[label=\Alph*.]
			\item Online store query latency and model inference time
			\item Training time
			\item Number of features in the training set only
			\item Cluster size
		\end{enumerate}
	\end{frame}
	\begin{frame}{Answer 334}
		\textbf{Correct: A} --- latency depends on online lookup and inference.\\
		\vspace{0.5em}
		\textbf{Wrong:}
		\begin{itemize}
			\item B, C, D are not the primary factors.
		\end{itemize}
	\end{frame}
	
	% ---------------------------------------------------------------
	\begin{frame}{Question 335}
		What is a batch prediction table?
		\begin{enumerate}[label=\Alph*.]
			\item A table storing offline model predictions
			\item A feature table
			\item A model
			\item A notebook
		\end{enumerate}
	\end{frame}
	\begin{frame}{Answer 335}
		\textbf{Correct: A} --- it stores offline predictions.\\
		\vspace{0.5em}
		\textbf{Wrong:}
		\begin{itemize}
			\item B, C, D are unrelated.
		\end{itemize}
	\end{frame}
	
	% ---------------------------------------------------------------
	\begin{frame}{Question 336}
		What is the purpose of a feature freshness policy?
		\begin{enumerate}[label=\Alph*.]
			\item To ensure features are updated within an acceptable window
			\item To train models
			\item To register models
			\item To deploy models
		\end{enumerate}
	\end{frame}
	\begin{frame}{Answer 336}
		\textbf{Correct: A} --- freshness policies bound feature staleness.\\
		\vspace{0.5em}
		\textbf{Wrong:}
		\begin{itemize}
			\item B, C, D are unrelated.
		\end{itemize}
	\end{frame}
	
	% ---------------------------------------------------------------
	\begin{frame}{Question 337}
		What is the main difference between a view and a table in Unity Catalog?
		\begin{enumerate}[label=\Alph*.]
			\item A view is a saved query; a table stores data
			\item A view stores data; a table is a query
			\item They are identical
			\item A view cannot be queried
		\end{enumerate}
	\end{frame}
	\begin{frame}{Answer 337}
		\textbf{Correct: A} --- views are saved queries.\\
		\vspace{0.5em}
		\textbf{Wrong:}
		\begin{itemize}
			\item B is reversed.
			\item C, D are false.
		\end{itemize}
	\end{frame}
	
	% ---------------------------------------------------------------
	\begin{frame}{Question 338}
		What is a managed table in Unity Catalog?
		\begin{enumerate}[label=\Alph*.]
			\item A table whose data lifecycle is managed by UC
			\item A table stored outside UC
			\item A view
			\item A model
		\end{enumerate}
	\end{frame}
	\begin{frame}{Answer 338}
		\textbf{Correct: A} --- managed tables have UC-managed storage.\\
		\vspace{0.5em}
		\textbf{Wrong:}
		\begin{itemize}
			\item B is external.
			\item C, D are other objects.
		\end{itemize}
	\end{frame}
	
	% ---------------------------------------------------------------
	\begin{frame}{Question 339}
		What is an external table in Unity Catalog?
		\begin{enumerate}[label=\Alph*.]
			\item A table whose data lives at a user-specified storage path
			\item A table managed by UC
			\item A view
			\item A model
		\end{enumerate}
	\end{frame}
	\begin{frame}{Answer 339}
		\textbf{Correct: A} --- external tables use user storage.\\
		\vspace{0.5em}
		\textbf{Wrong:}
		\begin{itemize}
			\item B is managed.
			\item C, D are other objects.
		\end{itemize}
	\end{frame}
	
	% ---------------------------------------------------------------
	\begin{frame}{Question 340}
		Which Unity Catalog privilege allows executing a model?
		\begin{enumerate}[label=\Alph*.]
			\item \texttt{EXECUTE}
			\item \texttt{SELECT}
			\item \texttt{MODIFY}
			\item \texttt{READ}
		\end{enumerate}
	\end{frame}
	\begin{frame}{Answer 340}
		\textbf{Correct: A} --- \texttt{EXECUTE} allows model use.\\
		\vspace{0.5em}
		\textbf{Wrong:}
		\begin{itemize}
			\item B is table read.
			\item C is table write.
			\item D is not a UC privilege.
		\end{itemize}
	\end{frame}
	
	% ---------------------------------------------------------------
	\begin{frame}{Question 341}
		What does a model version's description field store?
		\begin{enumerate}[label=\Alph*.]
			\item Human-readable notes about the version
			\item The training data
			\item The model weights
			\item The cluster config
		\end{enumerate}
	\end{frame}
	\begin{frame}{Answer 341}
		\textbf{Correct: A} --- descriptions store notes.\\
		\vspace{0.5em}
		\textbf{Wrong:}
		\begin{itemize}
			\item B, C, D are stored elsewhere.
		\end{itemize}
	\end{frame}
	
	% ---------------------------------------------------------------
	\begin{frame}{Question 342}
		What does a model version's source run link provide?
		\begin{enumerate}[label=\Alph*.]
			\item Traceability back to the originating MLflow run
			\item The cluster config
			\item The serving endpoint
			\item The feature table
		\end{enumerate}
	\end{frame}
	\begin{frame}{Answer 342}
		\textbf{Correct: A} --- source run links provide traceability.\\
		\vspace{0.5em}
		\textbf{Wrong:}
		\begin{itemize}
			\item B, C, D are unrelated.
		\end{itemize}
	\end{frame}
	
	% ---------------------------------------------------------------
	\begin{frame}{Question 343}
		What is a model version's stage used for in the workspace registry?
		\begin{enumerate}[label=\Alph*.]
			\item Indicating lifecycle state such as Staging or Production
			\item Storing metrics
			\item Training the model
			\item Storing data
		\end{enumerate}
	\end{frame}
	\begin{frame}{Answer 343}
		\textbf{Correct: A} --- stages indicate lifecycle state.\\
		\vspace{0.5em}
		\textbf{Wrong:}
		\begin{itemize}
			\item B, C, D are unrelated.
		\end{itemize}
	\end{frame}
	
	% ---------------------------------------------------------------
	\begin{frame}{Question 344}
		Why do Unity Catalog model registries favor aliases over stages?
		\begin{enumerate}[label=\Alph*.]
			\item Aliases are more flexible and support multiple deployment references
			\item Stages are deprecated in all tools
			\item Aliases store data
			\item Aliases train models
		\end{enumerate}
	\end{frame}
	\begin{frame}{Answer 344}
		\textbf{Correct: A} --- aliases are flexible deployment references.\\
		\vspace{0.5em}
		\textbf{Wrong:}
		\begin{itemize}
			\item B, C, D are false.
		\end{itemize}
	\end{frame}
	
	% ---------------------------------------------------------------
	\begin{frame}{Question 345}
		What is a model version's archive action used for?
		\begin{enumerate}[label=\Alph*.]
			\item Retiring a version from active use
			\item Training the model
			\item Registering a new version
			\item Deploying the model
		\end{enumerate}
	\end{frame}
	\begin{frame}{Answer 345}
		\textbf{Correct: A} --- archiving retires versions.\\
		\vspace{0.5em}
		\textbf{Wrong:}
		\begin{itemize}
			\item B, C, D are other actions.
		\end{itemize}
	\end{frame}
	
	% ---------------------------------------------------------------
	\begin{frame}{Question 346}
		What is a model version's comment used for?
		\begin{enumerate}[label=\Alph*.]
			\item Collaborators leave notes about the version
			\item Storing data
			\item Training the model
			\item Serving the model
		\end{enumerate}
	\end{frame}
	\begin{frame}{Answer 346}
		\textbf{Correct: A} --- comments store collaborator notes.\\
		\vspace{0.5em}
		\textbf{Wrong:}
		\begin{itemize}
			\item B, C, D are unrelated.
		\end{itemize}
	\end{frame}
	
	% ---------------------------------------------------------------
	\begin{frame}{Question 347}
		What is a MLflow model version's tag used for?
		\begin{enumerate}[label=\Alph*.]
			\item Attaching key-value metadata
			\item Storing metrics
			\item Storing features
			\item Storing the model weights
		\end{enumerate}
	\end{frame}
	\begin{frame}{Answer 347}
		\textbf{Correct: A} --- tags attach metadata.\\
		\vspace{0.5em}
		\textbf{Wrong:}
		\begin{itemize}
			\item B, C, D are stored elsewhere.
		\end{itemize}
	\end{frame}
	
	% ---------------------------------------------------------------
	\begin{frame}{Question 348}
		Which MLflow API sets a model version tag?
		\begin{enumerate}[label=\Alph*.]
			\item \texttt{mlflow.set\_model\_version\_tag()}
			\item \texttt{mlflow.log\_tag()}
			\item \texttt{mlflow.set\_tag()}
			\item \texttt{mlflow.tag\_model()}
		\end{enumerate}
	\end{frame}
	\begin{frame}{Answer 348}
		\textbf{Correct: A} --- it sets a model version tag.\\
		\vspace{0.5em}
		\textbf{Wrong:}
		\begin{itemize}
			\item B, C set run tags.
			\item D is not the API.
		\end{itemize}
	\end{frame}
	
	% ---------------------------------------------------------------
	\begin{frame}{Question 349}
		What is a model version's alias change recorded as?
		\begin{enumerate}[label=\Alph*.]
			\item A registry activity event
			\item A training log
			\item A feature table
			\item A notebook
		\end{enumerate}
	\end{frame}
	\begin{frame}{Answer 349}
		\textbf{Correct: A} --- alias changes are registry activity.\\
		\vspace{0.5em}
		\textbf{Wrong:}
		\begin{itemize}
			\item B, C, D are unrelated.
		\end{itemize}
	\end{frame}
	
	% ---------------------------------------------------------------
	\begin{frame}{Question 350}
		What is a Databricks notebook's default language?
		\begin{enumerate}[label=\Alph*.]
			\item Python
			\item SQL
			\item Scala
			\item R
		\end{enumerate}
	\end{frame}
	\begin{frame}{Answer 350}
		\textbf{Correct: A} --- Python is the default.\\
		\vspace{0.5em}
		\textbf{Wrong:}
		\begin{itemize}
			\item B, C, D are other supported languages.
		\end{itemize}
	\end{frame}
	
	% ---------------------------------------------------------------
	\begin{frame}{Question 351}
		What is a magic command in a Databricks notebook?
		\begin{enumerate}[label=\Alph*.]
			\item A command that changes the notebook's behavior or language
			\item A model
			\item A table
			\item A cluster
		\end{enumerate}
	\end{frame}
	\begin{frame}{Answer 351}
		\textbf{Correct: A} --- magics alter notebook behavior.\\
		\vspace{0.5em}
		\textbf{Wrong:}
		\begin{itemize}
			\item B, C, D are unrelated.
		\end{itemize}
	\end{frame}
	
	% ---------------------------------------------------------------
	\begin{frame}{Question 352}
		Which magic command switches a cell to SQL?
		\begin{enumerate}[label=\Alph*.]
			\item \texttt{\%sql}
			\item \texttt{\%python}
			\item \texttt{\%scala}
			\item \texttt{\%md}
		\end{enumerate}
	\end{frame}
	\begin{frame}{Answer 352}
		\textbf{Correct: A} --- \texttt{\%sql} switches to SQL.\\
		\vspace{0.5em}
		\textbf{Wrong:}
		\begin{itemize}
			\item B, C, D switch to other languages or Markdown.
		\end{itemize}
	\end{frame}
	
	% ---------------------------------------------------------------
	\begin{frame}{Question 353}
		Which magic command switches a cell to Markdown?
		\begin{enumerate}[label=\Alph*.]
			\item \texttt{\%md}
			\item \texttt{\%sql}
			\item \texttt{\%python}
			\item \texttt{\%sh}
		\end{enumerate}
	\end{frame}
	\begin{frame}{Answer 353}
		\textbf{Correct: A} --- \texttt{\%md} switches to Markdown.\\
		\vspace{0.5em}
		\textbf{Wrong:}
		\begin{itemize}
			\item B, C, D switch to other languages or shell.
		\end{itemize}
	\end{frame}
	
	% ---------------------------------------------------------------
	\begin{frame}{Question 354}
		Which magic command runs shell commands?
		\begin{enumerate}[label=\Alph*.]
			\item \texttt{\%sh}
			\item \texttt{\%sql}
			\item \texttt{\%python}
			\item \texttt{\%md}
		\end{enumerate}
	\end{frame}
	\begin{frame}{Answer 354}
		\textbf{Correct: A} --- \texttt{\%sh} runs shell commands.\\
		\vspace{0.5em}
		\textbf{Wrong:}
		\begin{itemize}
			\item B, C, D are other magics.
		\end{itemize}
	\end{frame}
	
	% ---------------------------------------------------------------
	\begin{frame}{Question 355}
		What does \texttt{dbutils.fs.ls()} do?
		\begin{enumerate}[label=\Alph*.]
			\item Lists files in a filesystem path
			\item Lists clusters
			\item Lists models
			\item Lists tables
		\end{enumerate}
	\end{frame}
	\begin{frame}{Answer 355}
		\textbf{Correct: A} --- it lists files.\\
		\vspace{0.5em}
		\textbf{Wrong:}
		\begin{itemize}
			\item B, C, D are other utilities.
		\end{itemize}
	\end{frame}
	
	% ---------------------------------------------------------------
	\begin{frame}{Question 356}
		What does \texttt{dbutils.notebook.run()} do?
		\begin{enumerate}[label=\Alph*.]
			\item Runs another notebook and returns its exit value
			\item Trains a model
			\item Registers a model
			\item Lists files
		\end{enumerate}
	\end{frame}
	\begin{frame}{Answer 356}
		\textbf{Correct: A} --- it runs a notebook and returns its value.\\
		\vspace{0.5em}
		\textbf{Wrong:}
		\begin{itemize}
			\item B, C, D are unrelated.
		\end{itemize}
	\end{frame}
	
	% ---------------------------------------------------------------
	\begin{frame}{Question 357}
		What does \texttt{dbutils.notebook.exit()} do?
		\begin{enumerate}[label=\Alph*.]
			\item Exits a notebook with a value
			\item Exits a cluster
			\item Exits a model
			\item Exits a table
		\end{enumerate}
	\end{frame}
	\begin{frame}{Answer 357}
		\textbf{Correct: A} --- it exits a notebook with a value.\\
		\vspace{0.5em}
		\textbf{Wrong:}
		\begin{itemize}
			\item B, C, D are unrelated.
		\end{itemize}
	\end{frame}
	
	% ---------------------------------------------------------------
	\begin{frame}{Question 358}
		What is a notebook task in a Databricks Job?
		\begin{enumerate}[label=\Alph*.]
			\item A single notebook executed as part of the job
			\item A cluster
			\item A model
			\item A table
		\end{enumerate}
	\end{frame}
	\begin{frame}{Answer 358}
		\textbf{Correct: A} --- a task runs a notebook.\\
		\vspace{0.5em}
		\textbf{Wrong:}
		\begin{itemize}
			\item B, C, D are unrelated.
		\end{itemize}
	\end{frame}
	
	% ---------------------------------------------------------------
	\begin{frame}{Question 359}
		What is a job dependency?
		\begin{enumerate}[label=\Alph*.]
			\item A task that runs after another task succeeds
			\item A cluster
			\item A model
			\item A table
		\end{enumerate}
	\end{frame}
	\begin{frame}{Answer 359}
		\textbf{Correct: A} --- dependencies order tasks.\\
		\vspace{0.5em}
		\textbf{Wrong:}
		\begin{itemize}
			\item B, C, D are unrelated.
		\end{itemize}
	\end{frame}
	
	% ---------------------------------------------------------------
	\begin{frame}{Question 360}
		What is a job trigger?
		\begin{enumerate}[label=\Alph*.]
			\item A schedule or event that starts a job
			\item A model
			\item A table
			\item A cluster
		\end{enumerate}
	\end{frame}
	\begin{frame}{Answer 360}
		\textbf{Correct: A} --- triggers start jobs.\\
		\vspace{0.5em}
		\textbf{Wrong:}
		\begin{itemize}
			\item B, C, D are unrelated.
		\end{itemize}
	\end{frame}
	
	% ---------------------------------------------------------------
	\begin{frame}{Question 361}
		What does a job retry policy control?
		\begin{enumerate}[label=\Alph*.]
			\item How many times and how often failed tasks are retried
			\item The model type
			\item The feature table
			\item The cluster size
		\end{enumerate}
	\end{frame}
	\begin{frame}{Answer 361}
		\textbf{Correct: A} --- retry policies control retries.\\
		\vspace{0.5em}
		\textbf{Wrong:}
		\begin{itemize}
			\item B, C, D are unrelated.
		\end{itemize}
	\end{frame}
	
	% ---------------------------------------------------------------
	\begin{frame}{Question 362}
		What is a cluster autoscaling minimum?
		\begin{enumerate}[label=\Alph*.]
			\item The smallest number of workers a cluster can scale to
			\item The runtime version
			\item The driver size
			\item The storage size
		\end{enumerate}
	\end{frame}
	\begin{frame}{Answer 362}
		\textbf{Correct: A} --- minimum workers bound autoscaling.\\
		\vspace{0.5em}
		\textbf{Wrong:}
		\begin{itemize}
			\item B, C, D are other settings.
		\end{itemize}
	\end{frame}
	
	% ---------------------------------------------------------------
	\begin{frame}{Question 363}
		What is a cluster autoscaling maximum?
		\begin{enumerate}[label=\Alph*.]
			\item The largest number of workers a cluster can scale to
			\item The runtime version
			\item The driver size
			\item The storage size
		\end{enumerate}
	\end{frame}
	\begin{frame}{Answer 363}
		\textbf{Correct: A} --- maximum workers bound autoscaling.\\
		\vspace{0.5em}
		\textbf{Wrong:}
		\begin{itemize}
			\item B, C, D are other settings.
		\end{itemize}
	\end{frame}
	
	% ---------------------------------------------------------------
	\begin{frame}{Question 364}
		What is a spot instance in Databricks?
		\begin{enumerate}[label=\Alph*.]
			\item A lower-cost, interruptible cloud instance
			\item A type of model
			\item A feature table
			\item A cluster policy
		\end{enumerate}
	\end{frame}
	\begin{frame}{Answer 364}
		\textbf{Correct: A} --- spot instances are cheaper but interruptible.\\
		\vspace{0.5em}
		\textbf{Wrong:}
		\begin{itemize}
			\item B, C, D are unrelated.
		\end{itemize}
	\end{frame}
	
	% ---------------------------------------------------------------
	\begin{frame}{Question 365}
		Why might spot instances be risky for long training jobs?
		\begin{enumerate}[label=\Alph*.]
			\item They can be reclaimed, interrupting the job
			\item They are slower
			\item They cannot run Python
			\item They cannot use GPUs
		\end{enumerate}
	\end{frame}
	\begin{frame}{Answer 365}
		\textbf{Correct: A} --- spot instances can be interrupted.\\
		\vspace{0.5em}
		\textbf{Wrong:}
		\begin{itemize}
			\item B, C, D are false.
		\end{itemize}
	\end{frame}
	
	% ---------------------------------------------------------------
	\begin{frame}{Question 366}
		What is a GPU cluster used for?
		\begin{enumerate}[label=\Alph*.]
			\item Deep learning and other compute-intensive workloads
			\item SQL BI only
			\item File storage
			\item Model serving only
		\end{enumerate}
	\end{frame}
	\begin{frame}{Answer 366}
		\textbf{Correct: A} --- GPUs accelerate heavy workloads.\\
		\vspace{0.5em}
		\textbf{Wrong:}
		\begin{itemize}
			\item B, C, D are not their main use.
		\end{itemize}
	\end{frame}
	
	% ---------------------------------------------------------------
	\begin{frame}{Question 367}
		What is the purpose of a cluster init script?
		\begin{enumerate}[label=\Alph*.]
			\item To run custom setup when a cluster starts
			\item To train a model
			\item To register a model
			\item To deploy a model
		\end{enumerate}
	\end{frame}
	\begin{frame}{Answer 367}
		\textbf{Correct: A} --- init scripts customize cluster startup.\\
		\vspace{0.5em}
		\textbf{Wrong:}
		\begin{itemize}
			\item B, C, D are unrelated.
		\end{itemize}
	\end{frame}
	
	% ---------------------------------------------------------------
	\begin{frame}{Question 368}
		What is the risk of long-running init scripts?
		\begin{enumerate}[label=\Alph*.]
			\item They slow cluster startup and can cause timeouts
			\item They improve accuracy
			\item They reduce storage
			\item They register models
		\end{enumerate}
	\end{frame}
	\begin{frame}{Answer 368}
		\textbf{Correct: A} --- slow scripts delay startup.\\
		\vspace{0.5em}
		\textbf{Wrong:}
		\begin{itemize}
			\item B, C, D are false.
		\end{itemize}
	\end{frame}
	
	% ---------------------------------------------------------------
	\begin{frame}{Question 369}
		What is the purpose of environment variables on a cluster?
		\begin{enumerate}[label=\Alph*.]
			\item To configure runtime behavior and credentials
			\item To train models
			\item To register models
			\item To store tables
		\end{enumerate}
	\end{frame}
	\begin{frame}{Answer 369}
		\textbf{Correct: A} --- env vars configure runtime behavior.\\
		\vspace{0.5em}
		\textbf{Wrong:}
		\begin{itemize}
			\item B, C, D are unrelated.
		\end{itemize}
	\end{frame}
	
	% ---------------------------------------------------------------
	\begin{frame}{Question 370}
		What does a notebook-scoped library installation do?
		\begin{enumerate}[label=\Alph*.]
			\item Installs a library only for the current notebook session
			\item Installs a library cluster-wide
			\item Installs a model
			\item Installs a table
		\end{enumerate}
	\end{frame}
	\begin{frame}{Answer 370}
		\textbf{Correct: A} --- notebook-scoped libraries are session-only.\\
		\vspace{0.5em}
		\textbf{Wrong:}
		\begin{itemize}
			\item B is cluster-scoped.
			\item C, D are false.
		\end{itemize}
	\end{frame}
	
	% ---------------------------------------------------------------
	\begin{frame}{Question 371}
		What does a cluster-scoped library installation do?
		\begin{enumerate}[label=\Alph*.]
			\item Installs a library for all notebooks on the cluster
			\item Installs a library for one notebook only
			\item Installs a model
			\item Installs a table
		\end{enumerate}
	\end{frame}
	\begin{frame}{Answer 371}
		\textbf{Correct: A} --- cluster-scoped libraries apply cluster-wide.\\
		\vspace{0.5em}
		\textbf{Wrong:}
		\begin{itemize}
			\item B is notebook-scoped.
			\item C, D are false.
		\end{itemize}
	\end{frame}
	
	% ---------------------------------------------------------------
	\begin{frame}{Question 372}
		What is a time-series forecasting problem?
		\begin{enumerate}[label=\Alph*.]
			\item Predicting future values from ordered historical data
			\item Clustering customers
			\item Classifying text
			\item Reducing dimensions
		\end{enumerate}
	\end{frame}
	\begin{frame}{Answer 372}
		\textbf{Correct: A} --- forecasting predicts future values.\\
		\vspace{0.5em}
		\textbf{Wrong:}
		\begin{itemize}
			\item B, C, D are other tasks.
		\end{itemize}
	\end{frame}
	
	% ---------------------------------------------------------------
	\begin{frame}{Question 373}
		What is the role of a lag feature in time-series forecasting?
		\begin{enumerate}[label=\Alph*.]
			\item It provides past values as predictors
			\item It encodes categories
			\item It reduces dimensions
			\item It scales features
		\end{enumerate}
	\end{frame}
	\begin{frame}{Answer 373}
		\textbf{Correct: A} --- lag features use past values.\\
		\vspace{0.5em}
		\textbf{Wrong:}
		\begin{itemize}
			\item B, C, D are unrelated.
		\end{itemize}
	\end{frame}
	
	% ---------------------------------------------------------------
	\begin{frame}{Question 374}
		What is a rolling window feature?
		\begin{enumerate}[label=\Alph*.]
			\item An aggregate over a moving window of past data
			\item A category encoding
			\item A dimensionality reduction
			\item A scaling method
		\end{enumerate}
	\end{frame}
	\begin{frame}{Answer 374}
		\textbf{Correct: A} --- rolling windows aggregate past data.\\
		\vspace{0.5em}
		\textbf{Wrong:}
		\begin{itemize}
			\item B, C, D are unrelated.
		\end{itemize}
	\end{frame}
	
	% ---------------------------------------------------------------
	\begin{frame}{Question 375}
		What is the risk of using future data in time-series features?
		\begin{enumerate}[label=\Alph*.]
			\item Data leakage and unrealistic performance
			\item Faster training
			\item Higher accuracy
			\item Smaller models
		\end{enumerate}
	\end{frame}
	\begin{frame}{Answer 375}
		\textbf{Correct: A} --- future data causes leakage.\\
		\vspace{0.5em}
		\textbf{Wrong:}
		\begin{itemize}
			\item B, C, D are false.
		\end{itemize}
	\end{frame}
	
	% ---------------------------------------------------------------
	\begin{frame}{Question 376}
		What is a recommendation system?
		\begin{enumerate}[label=\Alph*.]
			\item A model that suggests items to users
			\item A clustering model
			\item A forecasting model
			\item A dimensionality reduction model
		\end{enumerate}
	\end{frame}
	\begin{frame}{Answer 376}
		\textbf{Correct: A} --- recommenders suggest items.\\
		\vspace{0.5em}
		\textbf{Wrong:}
		\begin{itemize}
			\item B, C, D are other tasks.
		\end{itemize}
	\end{frame}
	
	% ---------------------------------------------------------------
	\begin{frame}{Question 377}
		Which Spark ML algorithm is used for collaborative filtering?
		\begin{enumerate}[label=\Alph*.]
			\item ALS
			\item KMeans
			\item LinearRegression
			\item RandomForest
		\end{enumerate}
	\end{frame}
	\begin{frame}{Answer 377}
		\textbf{Correct: A} --- ALS is for collaborative filtering.\\
		\vspace{0.5em}
		\textbf{Wrong:}
		\begin{itemize}
			\item B, C, D are other algorithms.
		\end{itemize}
	\end{frame}
	
	% ---------------------------------------------------------------
	\begin{frame}{Question 378}
		What does a graph analytics library help with?
		\begin{enumerate}[label=\Alph*.]
			\item Analyzing relationships between entities
			\item Training image models
			\item Encoding text
			\item Scaling features
		\end{enumerate}
	\end{frame}
	\begin{frame}{Answer 378}
		\textbf{Correct: A} --- graph analytics studies relationships.\\
		\vspace{0.5em}
		\textbf{Wrong:}
		\begin{itemize}
			\item B, C, D are unrelated.
		\end{itemize}
	\end{frame}
	
	% ---------------------------------------------------------------
	\begin{frame}{Question 379}
		What is a graph in graph analytics?
		\begin{enumerate}[label=\Alph*.]
			\item A set of nodes and edges
			\item A plot
			\item A table
			\item A model
		\end{enumerate}
	\end{frame}
	\begin{frame}{Answer 379}
		\textbf{Correct: A} --- graphs have nodes and edges.\\
		\vspace{0.5em}
		\textbf{Wrong:}
		\begin{itemize}
			\item B, C, D are not graphs.
		\end{itemize}
	\end{frame}
	
	% ---------------------------------------------------------------
	\begin{frame}{Question 380}
		What is a unit test in ML code?
		\begin{enumerate}[label=\Alph*.]
			\item A test of a small piece of logic
			\item A full training run
			\item A model
			\item A cluster
		\end{enumerate}
	\end{frame}
	\begin{frame}{Answer 380}
		\textbf{Correct: A} --- unit tests check small logic pieces.\\
		\vspace{0.5em}
		\textbf{Wrong:}
		\begin{itemize}
			\item B, C, D are not unit tests.
		\end{itemize}
	\end{frame}
	
	% ---------------------------------------------------------------
	\begin{frame}{Question 381}
		What is a data validation test in an ML pipeline?
		\begin{enumerate}[label=\Alph*.]
			\item A check that data meets expected schema and quality rules
			\item A model evaluation
			\item A cluster test
			\item A serving test
		\end{enumerate}
	\end{frame}
	\begin{frame}{Answer 381}
		\textbf{Correct: A} --- data validation enforces schema and quality.\\
		\vspace{0.5em}
		\textbf{Wrong:}
		\begin{itemize}
			\item B, C, D are other tests.
		\end{itemize}
	\end{frame}
	
	% ---------------------------------------------------------------
	\begin{frame}{Question 382}
		What is a model validation test?
		\begin{enumerate}[label=\Alph*.]
			\item A check that model performance meets thresholds
			\item A data schema check
			\item A cluster check
			\item A serving check
		\end{enumerate}
	\end{frame}
	\begin{frame}{Answer 382}
		\textbf{Correct: A} --- model validation checks performance.\\
		\vspace{0.5em}
		\textbf{Wrong:}
		\begin{itemize}
			\item B, C, D are other checks.
		\end{itemize}
	\end{frame}
	
	% ---------------------------------------------------------------
	\begin{frame}{Question 383}
		Why is CI/CD important for ML?
		\begin{enumerate}[label=\Alph*.]
			\item It automates testing and deployment of code and models
			\item It trains faster
			\item It stores data
			\item It reduces feature count
		\end{enumerate}
	\end{frame}
	\begin{frame}{Answer 383}
		\textbf{Correct: A} --- CI/CD automates testing and deployment.\\
		\vspace{0.5em}
		\textbf{Wrong:}
		\begin{itemize}
			\item B, C, D are false.
		\end{itemize}
	\end{frame}
	
	% ---------------------------------------------------------------
	\begin{frame}{Question 384}
		What is a model deployment pipeline?
		\begin{enumerate}[label=\Alph*.]
			\item An automated sequence that validates and deploys models
			\item A training notebook
			\item A feature table
			\item A cluster
		\end{enumerate}
	\end{frame}
	\begin{frame}{Answer 384}
		\textbf{Correct: A} --- deployment pipelines automate release.\\
		\vspace{0.5em}
		\textbf{Wrong:}
		\begin{itemize}
			\item B, C, D are unrelated.
		\end{itemize}
	\end{frame}
	
	% ---------------------------------------------------------------
	\begin{frame}{Question 385}
		What is a rollback in model deployment?
		\begin{enumerate}[label=\Alph*.]
			\item Reverting to a previous model version
			\item Training a new model
			\item Registering a model
			\item Deleting data
		\end{enumerate}
	\end{frame}
	\begin{frame}{Answer 385}
		\textbf{Correct: A} --- rollback reverts to an earlier version.\\
		\vspace{0.5em}
		\textbf{Wrong:}
		\begin{itemize}
			\item B, C, D are unrelated.
		\end{itemize}
	\end{frame}
	
	% ---------------------------------------------------------------
	\begin{frame}{Question 386}
		What is a model alias in Unity Catalog used for during deployment?
		\begin{enumerate}[label=\Alph*.]
			\item It points to the version served by an endpoint
			\item It stores data
			\item It trains the model
			\item It deletes the model
		\end{enumerate}
	\end{frame}
	\begin{frame}{Answer 386}
		\textbf{Correct: A} --- aliases point endpoints to versions.\\
		\vspace{0.5em}
		\textbf{Wrong:}
		\begin{itemize}
			\item B, C, D are false.
		\end{itemize}
	\end{frame}
	
	% ---------------------------------------------------------------
	\begin{frame}{Question 387}
		What is A/B testing in model deployment?
		\begin{enumerate}[label=\Alph*.]
			\item Comparing performance of two model versions in production
			\item Training two models
			\item Deleting two models
			\item Registering two models
		\end{enumerate}
	\end{frame}
	\begin{frame}{Answer 387}
		\textbf{Correct: A} --- A/B testing compares deployed versions.\\
		\vspace{0.5em}
		\textbf{Wrong:}
		\begin{itemize}
			\item B, C, D are not A/B testing.
		\end{itemize}
	\end{frame}
	
	% ---------------------------------------------------------------
	\begin{frame}{Question 388}
		How is traffic splitting used for A/B testing?
		\begin{enumerate}[label=\Alph*.]
			\item It routes a fraction of requests to each model version
			\item It trains models
			\item It deletes models
			\item It registers models
		\end{enumerate}
	\end{frame}
	\begin{frame}{Answer 388}
		\textbf{Correct: A} --- traffic splitting routes requests.\\
		\vspace{0.5em}
		\textbf{Wrong:}
		\begin{itemize}
			\item B, C, D are false.
		\end{itemize}
	\end{frame}
	
	% ---------------------------------------------------------------
	\begin{frame}{Question 389}
		What is a model's inference latency?
		\begin{enumerate}[label=\Alph*.]
			\item The time to produce a prediction
			\item The training time
			\item The storage size
			\item The feature count
		\end{enumerate}
	\end{frame}
	\begin{frame}{Answer 389}
		\textbf{Correct: A} --- latency is prediction time.\\
		\vspace{0.5em}
		\textbf{Wrong:}
		\begin{itemize}
			\item B, C, D are unrelated.
		\end{itemize}
	\end{frame}
	
	% ---------------------------------------------------------------
	\begin{frame}{Question 390}
		What is throughput in model serving?
		\begin{enumerate}[label=\Alph*.]
			\item The number of predictions served per unit time
			\item The training time
			\item The storage size
			\item The feature count
		\end{enumerate}
	\end{frame}
	\begin{frame}{Answer 390}
		\textbf{Correct: A} --- throughput is predictions per time.\\
		\vspace{0.5em}
		\textbf{Wrong:}
		\begin{itemize}
			\item B, C, D are unrelated.
		\end{itemize}
	\end{frame}
	
	% ---------------------------------------------------------------
	\begin{frame}{Question 391}
		Why might batch inference be preferred over real-time serving?
		\begin{enumerate}[label=\Alph*.]
			\item Predictions are not needed immediately and batch is cheaper
			\item Batch is lower latency
			\item Batch requires a REST API
			\item Batch cannot use models
		\end{enumerate}
	\end{frame}
	\begin{frame}{Answer 391}
		\textbf{Correct: A} --- batch suits non-urgent scoring.\\
		\vspace{0.5em}
		\textbf{Wrong:}
		\begin{itemize}
			\item B, C, D are false.
		\end{itemize}
	\end{frame}
	
	% ---------------------------------------------------------------
	\begin{frame}{Question 392}
		What is the role of a feature table in batch inference?
		\begin{enumerate}[label=\Alph*.]
			\item It stores features joined to the scoring data
			\item It trains the model
			\item It serves the model
			\item It registers the model
		\end{enumerate}
	\end{frame}
	\begin{frame}{Answer 392}
		\textbf{Correct: A} --- feature tables supply features.\\
		\vspace{0.5em}
		\textbf{Wrong:}
		\begin{itemize}
			\item B, C, D are unrelated.
		\end{itemize}
	\end{frame}
	
	% ---------------------------------------------------------------
	\begin{frame}{Question 393}
		What is a model's input schema in serving?
		\begin{enumerate}[label=\Alph*.]
			\item The expected structure of incoming data
			\item The model weights
			\item The training time
			\item The cluster size
		\end{enumerate}
	\end{frame}
	\begin{frame}{Answer 393}
		\textbf{Correct: A} --- input schema defines expected data.\\
		\vspace{0.5em}
		\textbf{Wrong:}
		\begin{itemize}
			\item B, C, D are unrelated.
		\end{itemize}
	\end{frame}
	
	% ---------------------------------------------------------------
	\begin{frame}{Question 394}
		What happens if serving input violates the model signature?
		\begin{enumerate}[label=\Alph*.]
			\item The request may fail validation
			\item The model retrains
			\item The endpoint deletes itself
			\item The data is stored
		\end{enumerate}
	\end{frame}
	\begin{frame}{Answer 394}
		\textbf{Correct: A} --- signature validation can reject input.\\
		\vspace{0.5em}
		\textbf{Wrong:}
		\begin{itemize}
			\item B, C, D are false.
		\end{itemize}
	\end{frame}
	
	% ---------------------------------------------------------------
	\begin{frame}{Question 395}
		What is a model's output schema in serving?
		\begin{enumerate}[label=\Alph*.]
			\item The structure of predictions returned by the model
			\item The model weights
			\item The training time
			\item The cluster size
		\end{enumerate}
	\end{frame}
	\begin{frame}{Answer 395}
		\textbf{Correct: A} --- output schema defines returned predictions.\\
		\vspace{0.5em}
		\textbf{Wrong:}
		\begin{itemize}
			\item B, C, D are unrelated.
		\end{itemize}
	\end{frame}
	
	% ---------------------------------------------------------------
	\begin{frame}{Question 396}
		What is the benefit of a feature store for real-time serving?
		\begin{enumerate}[label=\Alph*.]
			\item It provides low-latency online feature lookups
			\item It trains models
			\item It registers models
			\item It stores notebooks
		\end{enumerate}
	\end{frame}
	\begin{frame}{Answer 396}
		\textbf{Correct: A} --- online tables serve low-latency features.\\
		\vspace{0.5em}
		\textbf{Wrong:}
		\begin{itemize}
			\item B, C, D are unrelated.
		\end{itemize}
	\end{frame}
	
	% ---------------------------------------------------------------
	\begin{frame}{Question 397}
		What is a common cause of training-serving skew?
		\begin{enumerate}[label=\Alph*.]
			\item Features computed differently in training and serving
			\item Using too many workers
			\item Using too few features
			\item Using a large model
		\end{enumerate}
	\end{frame}
	\begin{frame}{Answer 397}
		\textbf{Correct: A} --- inconsistent feature logic causes skew.\\
		\vspace{0.5em}
		\textbf{Wrong:}
		\begin{itemize}
			\item B, C, D are not the cause.
		\end{itemize}
	\end{frame}
	
	% ---------------------------------------------------------------
	\begin{frame}{Question 398}
		How can Feature Store reduce training-serving skew?
		\begin{enumerate}[label=\Alph*.]
			\item By logging feature metadata with the model for consistent lookups
			\item By training more models
			\item By deleting features
			\item By reducing cluster size
		\end{enumerate}
	\end{frame}
	\begin{frame}{Answer 398}
		\textbf{Correct: A} --- logged metadata enables consistent feature lookups.\\
		\vspace{0.5em}
		\textbf{Wrong:}
		\begin{itemize}
			\item B, C, D do not reduce skew.
		\end{itemize}
	\end{frame}
	
	% ---------------------------------------------------------------
	\begin{frame}{Question 399}
		Why is monitoring prediction distributions important?
		\begin{enumerate}[label=\Alph*.]
			\item Shifts can indicate drift or data issues
			\item It trains models
			\item It registers models
			\item It deletes models
		\end{enumerate}
	\end{frame}
	\begin{frame}{Answer 399}
		\textbf{Correct: A} --- prediction shifts signal drift.\\
		\vspace{0.5em}
		\textbf{Wrong:}
		\begin{itemize}
			\item B, C, D are unrelated.
		\end{itemize}
	\end{frame}
	
	% ---------------------------------------------------------------
	\begin{frame}{Question 400}
		What is the best practice for retiring a model version in production?
		\begin{enumerate}[label=\Alph*.]
			\item Move the serving alias to a new version before archiving the old one
			\item Delete the model immediately
			\item Retrain from scratch
			\item Stop all monitoring
		\end{enumerate}
	\end{frame}
	\begin{frame}{Answer 400}
		\textbf{Correct: A} --- move the alias first, then archive.\\
		\vspace{0.5em}
		\textbf{Wrong:}
		\begin{itemize}
			\item B causes downtime.
			\item C is unnecessary.
			\item D removes safeguards.
		\end{itemize}
	\end{frame}
	
\end{document}